%=============================================================
%  Module I.4 / L04 — Rabi Model and Two-Level Systems
%  Series I > Module I.4 > Lecture L04
%  Duration: 90 minutes  |  All tiers (T1–T5)
%=============================================================
\renewcommand{\lectureCode}{I.4 / L04}

\section*{L04 — Rabi Model and Two-Level Systems}
\label{sec:L04}
\addcontentsline{toc}{section}{L04 — Rabi Model and Two-Level Systems}

%--- Lecture header
\begin{tcolorbox}[lecturebox, title={Module I.4 / L04 — Metadata}]
\begin{tabular}{@{}ll@{}}
\textbf{Series:}    & I — Introductory Quantum Mechanics \\
\textbf{Module:}    & I.4 — Paradigm Physical Systems \\
\textbf{Lecture:}   & L04 — Rabi Model and Two-Level Systems \\
\textbf{Duration:}  & 90 minutes \\
\textbf{Tiers:}     & T1 (HS) · T2 (BegUG) · T3 (AdvUG) · T4 (MSc) · T5 (PhD) \\
\textbf{Preceding:} & I.4 / L03 — Harmonic Oscillator \\
\textbf{Following:} & I.4 / L05 — Minimal Coupling and Landau Levels \\
\textbf{Prerequisites:} & Modules I.1–I.3 complete \\
\end{tabular}

\medskip
\textbf{Key concepts:} TLS Bloch sphere, rotating wave approximation, Rabi oscillations, Jaynes–Cummings model
\end{tcolorbox}

%------------------------------------------------------------
\subsection*{L04.0 — Learning Objectives (per tier)}
\label{sec:L04.0}

\tiertag{tier1col}{Tier 1 — HS}\quad
Identify the main physical phenomenon studied in this lecture; describe it in
non-mathematical terms; state one real-world application.

\tiertag{tier2col}{Tier 2 — BegUG}\quad
Apply the key equations to compute numerical results; verify consistency with
known limits; translate between physical and mathematical descriptions.

\tiertag{tier3col}{Tier 3 — AdvUG}\quad
Derive central results from first principles; prove key identities; identify
the mathematical structure underlying the physical phenomenon.

\tiertag{tier4col}{Tier 4 — MSc}\quad
Analyse formal extensions; connect to symmetry principles; generalise to
related systems; identify the role of this lecture in the broader programme.

\tiertag{tier5col}{Tier 5 — PhD}\quad
Establish rigorous mathematical foundations; identify open research problems;
connect to modern literature; critique standard textbook treatments.

%--- Lecture content -----------------------------------------

\subsection*{L04.1 — Two-Level System Hamiltonian}
\begin{tcolorbox}[lecturebox, title={TLS and Bloch Sphere}]
$\hat{H}_0 = \frac{\hbar\omega_0}{2}\hat{\sigma}_z$; general state:
$\ket{\psi}=\cos(\theta/2)\ket{\uparrow}+e^{i\phi}\sin(\theta/2)\ket{\downarrow}$
(point on Bloch sphere).
\end{tcolorbox}

\subsection*{L04.2 — Driven TLS and Rabi Oscillations}
In the rotating frame with RWA, the Hamiltonian becomes:
\begin{equation}
  \hat{H}_{\rm rot} = \frac{\hbar\Delta}{2}\hat{\sigma}_z + \frac{\hbar\Omega_R}{2}\hat{\sigma}_x,
  \quad \Delta = \omega - \omega_0,\quad \Omega_R = dE_0/\hbar.
  \label{eq:L04.1}
\end{equation}
Population inversion: $P_\uparrow(t) = (\Omega_R/\tilde\Omega)^2\sin^2(\tilde\Omega t/2)$,
$\tilde\Omega = \sqrt{\Delta^2+\Omega_R^2}$.

\subsection*{L04.3 — Jaynes–Cummings Model}
Full quantum treatment: TLS + single photon mode:
\begin{equation}
  \hat{H}_{\rm JC} = \hbar\omega_c\hat{a}^\dagger\hat{a}
    + \tfrac{\hbar\omega_0}{2}\hat{\sigma}_z
    + \hbar g(\hat{a}\hat{\sigma}_+ + \hat{a}^\dagger\hat{\sigma}_-).
  \label{eq:L04.2}
\end{equation}
Dressed states; collapse and revival of Rabi oscillations for coherent photon states.


%------------------------------------------------------------
\subsection*{L04.C — Connections}
\label{sec:L04.C}
\textbf{Backward:} I.4/L03 — Harmonic Oscillator and Modules I.1–I.3.
\textbf{Forward:} I.4/L05 — Minimal Coupling and Landau Levels builds directly on the formalism developed here.
\textbf{Cross-module:} Series~II modules extend these results to many-body and
advanced settings; Series~III applies them to molecular electronic structure.

%=============================================================
%  ASSESSMENT BUNDLE — L04
%=============================================================
\newpage
\section*{L04 — Assessment Artifact Bundle}

\begin{tcolorbox}[ugconceptbox, title={SCQU · L04 · Tiers 1–3 · 15–20 min}]
\textit{Ten simple concept questions (mix MC/TF/short) targeting the core results
of Rabi Model and Two-Level Systems. See artifact file \texttt{L04_artifacts.tex} for full questions.}

\textbf{Rubric:} 2 pts each (1.5 correct + 0.5 justification). Total: 20 pts.
\end{tcolorbox}

\begin{tcolorbox}[ugconceptbox, title={CCQU · L04 · Tiers 2–3 · 25–35 min}]
\textit{Ten complex concept questions: ``explain why,'' ``what would change if,''
``identify the flaw.'' Full questions in \texttt{L04_artifacts.tex}.}

\textbf{Rubric:} 2 pts each (1 key point + 1 argument). Total: 20 pts.
\end{tcolorbox}

\begin{tcolorbox}[ugprobbox, title={SPSU · L04 · Tiers 1–3 · 60–90 min}]
\textit{Ten problems (Part A: mechanics, Part B: interpretation, Part C: translation).
Full problems in \texttt{L04_artifacts.tex}.}

\textbf{Rubric:} Result 60\%, Method 30\%, Interpretation 10\%.
\end{tcolorbox}

\begin{tcolorbox}[ugprobbox, title={CPSU · L04 · Tier 3 · 3–6 hrs (4-layer)}]
\textit{Five multi-part derivation problems with four-layer scaffolding.
Layers 1–4 in \texttt{L04_ex01\_problem.tex}, \texttt{hints}, \texttt{adv\_hints}, \texttt{solution}.}

\textbf{Rubric:} Assumptions 15\%, Proof structure 45\%, Algebra 25\%, Insight 15\%.
\end{tcolorbox}

\begin{tcolorbox}[ugprobbox, title={SPJU + CPJU · L04 · Tiers 2–3}]
\textit{Simple project (2–6 hrs): structured computational/analytical task.
Complex project (8–15 hrs): open-ended synthesis.
Full specs in \texttt{L04_artifacts.tex}.}
\end{tcolorbox}

\begin{tcolorbox}[gradconceptbox, title={SCQG + CCQG · L04 · Tiers 4–5}]
\textit{Ten simple + ten complex concept questions at graduate level (rigour, domain
conditions, named theorems). Full questions in \texttt{L04_artifacts.tex}.}
\end{tcolorbox}

\begin{tcolorbox}[gradprobbox, title={SPSG + CPSG · L04 · Tiers 4–5}]
\textit{Ten simple + five complex graduate problems.
CPSG with four-layer format. See \texttt{L04_artifacts.tex}.}
\end{tcolorbox}

\begin{tcolorbox}[gradprobbox, title={SPJG (×5) + CPJG (×5) · L04 · Tiers 4–5}]
\textit{Five simple (3–8 hrs each) and five complex (10–20 hrs each) graduate projects.
Full specs in \texttt{L04_artifacts.tex}.}
\end{tcolorbox}

\begin{tcolorbox}[researchbox, title={WRQ (×5) + ORQ (×5) · L04 · Tiers 4–5}]
\textit{Five well-defined research questions (5–10 hrs each) and five open-ended
research questions (10–20 hrs each). Full prompts in \texttt{L04_artifacts.tex}.}
\end{tcolorbox}
